\section{Health: Care at a Distance}
\label{sec:health}

Health care is where the recurring tax becomes personal. A small, remote
health system pays it in four currencies: \emph{access} (specialist care,
advanced diagnostics and complex surgery often mean an overseas referral);
\emph{supply} (medicines, vaccines and equipment arrive in small, expensive,
infrequent consignments); \emph{diet} (imported, processed food displaces
local produce); and \emph{workforce} (doctors and nurses are scarce and
mobile).

\paragraph{Workforce.} \Cref{fig:doctors} plots physicians per 1,000 people
against remoteness. The median small island state has \docSIDSMedian{}
doctors per 1,000 and the median Pacific state \docPacificMedian, against
\docHICMedian{} in high-income economies. Remoteness and medical density are
correlated (\corrDocRemote), and the association survives a control for
income: comparing countries with similar income per head, a tenfold increase
in remoteness goes with \docRemoteCoef{} fewer physicians per 1,000 (standard
error \docRemoteSE). In a population of a hundred thousand, the absence of a
single specialist can mean an entire field of medicine is practised only by
referral abroad.

\begin{figure}[t]
\centering
\begin{tikzpicture}
\begin{axis}[edge, height=0.58\linewidth,
    xlabel={Economic remoteness, thousand km}, xmin=5.5, xmax=14.5,
    ymode=log, ymin=0.02, ymax=10, ytick={0.02,0.05,0.1,0.2,0.5,1,2,5,10},
    yticklabels={0.02,0.05,0.1,0.2,0.5,1,2,5,10},
    ylabel={Physicians per 1,000 people (log scale)},
    legend pos=south west]
  \addplot[dotsBase, only if={sids}{0}] table[x=remote_kkm,y=doctors] {\datadir/plot_countries.csv};
  \addlegendentry{Other economies}
  \addplot[dotsSIDS, only if={sids}{1}] table[x=remote_kkm,y=doctors] {\datadir/plot_countries.csv};
  \addlegendentry{Small island states}
  \addplot[labelled=lab_doc] table[x=remote_kkm,y=doctors] {\datadir/plot_countries.csv};
\end{axis}
\end{tikzpicture}
\caption{\takeaway{The further from the centre, the thinner the medical
  workforce.} Median physicians per 1,000: \docSIDSMedian{} in small island
  states, \docPacificMedian{} in the Pacific, \docHICMedian{} in high-income
  economies. Holding income constant, a tenfold rise in remoteness is associated with
  \docRemoteCoef{} fewer doctors per 1,000 (s.e.\ \docRemoteSE). Latest year 2013--23.
  \source{World Bank WDI (WHO data); CEPII; author's remoteness index.}}
\label{fig:doctors}
\end{figure}


\paragraph{Diet and disease.} The Pacific has some of the world's highest
rates of diabetes. In NCD-RisC's estimates \parencite{ncdrisc2024}, the
highest age-standardised prevalence in 2022 was in \diabTopName{}
(\diabTopValue\%), and \diabTopTenPacific{} of the ten highest-prevalence
countries are Pacific island states (\cref{fig:diabetes}). Imported, cheap,
energy-dense foods---refined flour, sugar, fatty meat cuts---have displaced
traditional diets, and Pacific governments have experimented with trade and
tax measures to push back \parencite{snowdon2013}.

\begin{figure}[t]
\centering
\begin{tikzpicture}
\begin{axis}[edge, height=0.52\linewidth,
    xmin=1990, xmax=2022, xtick={1990,2000,2010,2022},
    x tick label style={/pgf/number format/1000 sep={}},
    ymin=0, ymax=36, ylabel={Adults with diabetes, \% (age-standardised)},
    legend style={at={(0.02,0.98)},anchor=north west}]
  \addplot[lineEmber] table[x=year,y=pacific] {\datadir/diabetes_groups.csv};
  \addlegendentry{Pacific island states (median)}
  \addplot[lineTeal] table[x=year,y=sids] {\datadir/diabetes_groups.csv};
  \addlegendentry{All small island states (median)}
  \addplot[lineSlate] table[x=year,y=all] {\datadir/diabetes_groups.csv};
  \addlegendentry{All countries (median)}
\end{axis}
\end{tikzpicture}
\caption{\takeaway{Island diabetes rates were high and kept rising.}
  The median Pacific state went from \diabPacificNinety\% to \diabPacificNow\%
  between 1990 and 2022; the median country from \diabAllNinety\% to
  \diabAllNow\%. The rest of the world is catching up, so the relative gap
  narrowed (from \diabPacificRatioNinety{} to \diabPacificRatioNow{} times) even
  as the absolute gap held. Mean of male and female estimates, adults 18+.
  \source{NCD-RisC \parencite{ncdrisc2024}; UN M49.}}
\label{fig:diabetes}
\end{figure}


Two cautions apply. First, diet change has many causes: urbanisation,
incomes, physical activity, genetics, marketing and culture all play a part,
and the data here cannot apportion blame to remoteness or trade. The claim
is only that remote, import-dependent food systems make healthy diets
\emph{more expensive} and treatment harder to deliver---a tax on the cure as
well as a contributor to the disease. Second, the relative gap is narrowing,
for an unhappy reason: the rest of the world is catching up. The Pacific
median was \diabPacificRatioNinety{} times the global median in 1990 and
\diabPacificRatioNow{} times in 2022. What remote systems lack is not
awareness of the problem but the scale to manage it: dialysis units, retinal
screening, endocrinologists and reliable insulin supply chains.
